\section{Foundation Models and Learning at Scale}

Sections 4.1--4.3 argued that modern learning increasingly separates the acquisition of representation from the final downstream use of that representation.
Foundation models are a large-scale realization of that idea.
They are not best understood as ``very large neural networks'' in the abstract.
Rather, they are systems trained on broad data with a reusable interface, so that one parameter set can support many later tasks through prompting, fine-tuning, retrieval, or other adaptation mechanisms.
This section gives a sober account of that regime.
Its aim is not to celebrate scale, but to make precise what changes when representation learning, pretraining, and deployment are organized around one large pretrained model.

\subsection*{A precise working definition}

We begin with a definition tailored to the logic of this chapter.

\begin{definition}[Foundation model]
A \emph{foundation model} is a pretrained parametric model
\[
F_{\theta_0}
\]
obtained from broad data distribution $P_{\mathrm{pre}}$ by optimizing a generic pretraining objective
\[
\mathcal{L}_{\mathrm{pre}}(\theta),
\]
such that the resulting parameter vector $\theta_0$ supports a family of downstream tasks $\mathcal{T}$ through one or more adaptation mechanisms and deployment interfaces, with comparatively modest task-specific modification relative to training the model from scratch for each task.
\end{definition}

This definition is intentionally narrower than ``large model'' and broader than ``language model.''
Three ingredients matter.
First, the model is trained on broad rather than narrowly task-specific data.
Second, the pretraining objective is generic enough to produce reusable internal structure.
Third, the resulting system is reused across many tasks through a stable interface.
The word ``foundation'' therefore refers to a role in a learning pipeline, not merely to parameter count.

A small model can in principle play this role within a restricted domain, and a large model can fail to play it if it is trained only for one narrow benchmark.
In practice, however, the models usually discussed under this heading are large because broad reuse often benefits from large capacity, diverse data, and substantial compute.

\subsection*{Pretraining objective, adaptation mechanism, and deployment interface}

One source of confusion in current discussion is that several distinct layers are mixed together.
For a mathematically cleaner account, it is useful to separate them.

\begin{definition}[Pretraining objective]
A \emph{pretraining objective} is a loss functional
\[
\mathcal{L}_{\mathrm{pre}}(\theta)
=
\mathbb{E}_{x \sim P_{\mathrm{pre}}}\bigl[\ell_{\mathrm{pre}}(x;\theta)\bigr]
\]
or its empirical approximation on a large dataset.
Examples include masked prediction, autoregressive next-token prediction, denoising, contrastive objectives, or multimodal matching objectives.
Its role is to determine what regularities the model is rewarded for internalizing during large-scale training.
\end{definition}

\begin{definition}[Adaptation mechanism]
An \emph{adaptation mechanism} is a procedure that converts a pretrained model into a task-specialized system.
Formally, it is a map of the form
\[
\mathcal{A} : (\theta_0,\mathcal{D}_T) \mapsto \Theta_T,
\]
where $\mathcal{D}_T$ is task-specific information and $\Theta_T$ denotes either updated parameters, auxiliary modules, prompts, retrieval state, or some combination thereof.
Full fine-tuning, linear probing, adapters, low-rank adaptation, and instruction tuning are examples.
\end{definition}

\begin{definition}[Deployment interface]
A \emph{deployment interface} is the protocol by which a user or downstream system queries the pretrained or adapted model.
It specifies what object may be presented to the model at inference time and what output is requested.
For example, a prompt interface may accept text context $c$ and return a conditional distribution over continuations,
\[
y \mapsto p_{\theta}(y \mid c),
\]
while a classification interface may accept an input $x$ and return a label or score.
\end{definition}

These three notions should not be conflated.
The pretraining objective explains how the shared representation was learned.
The adaptation mechanism explains whether and how the parameters or auxiliary modules are changed for a new task.
The deployment interface explains how the resulting system is queried at use time.
Much confusion around foundation models arises when prompting is discussed as if it were a training objective, or instruction tuning is discussed as if it were only an inference trick.

\begin{example}[One model, three different layers]
Consider a decoder-only transformer trained with next-token prediction.
Its pretraining objective is autoregressive likelihood.
Its adaptation mechanism might be none at all (zero-shot prompting), or supervised instruction tuning, or low-rank fine-tuning.
Its deployment interface might be a text prompt followed by generation, a chat interface, or an API that returns log-probabilities over candidate outputs.
These are related, but they are not the same mathematical object.
\end{example}

\subsection*{A formal view of prompting and conditional behavior}

For concreteness, suppose a pretrained decoder defines
\[
p_{\theta}(x_{1:T})
=
\prod_{t=1}^{T} p_{\theta}(x_t \mid x_{<t}).
\]
If a context or prompt $c$ is provided, then the model induces a conditional distribution on continuations:
\[
p_{\theta}(y \mid c).
\]
This is the basic mathematical content of prompting.
One does not change the parameters.
One changes the conditioning context.

\begin{definition}[Prompting]
Let $\theta$ be fixed.
\emph{Prompting} is the use of a context variable $c$ at inference time to select a conditional prediction rule
\[
q_{\theta,c}(y) = p_{\theta}(y \mid c).
\]
The prompt may contain an instruction, demonstrations, formatting constraints, retrieved external information, or a question to be answered.
\end{definition}

This definition is deliberately simple.
It captures both ordinary continuation and more elaborate prompting patterns.
Few-shot prompting is the case in which the prompt includes example input--output pairs.
Chain-of-thought prompting is the case in which the prompt encourages intermediate reasoning-like text.
Retrieval-augmented prompting is the case in which the prompt is extended by externally fetched context.
All of these remain, mathematically, conditional inference with fixed parameters.

\begin{proposition}[Prompting changes context, not weights]
Fix model parameters $\theta$.
Then the family of predictors available through prompting is
\[
\mathcal{Q}_{\theta}
=
\{q_{\theta,c} : c \in \mathcal{C}\},
\qquad
q_{\theta,c}(y)=p_{\theta}(y \mid c).
\]
Prompting can select different conditional behaviors within $\mathcal{Q}_{\theta}$, but it does not alter the underlying parameter vector and therefore does not perform parameter adaptation.
\end{proposition}

\paragraph{Proof sketch.}
For fixed $\theta$, every admissible prompt $c$ defines a conditional distribution $p_{\theta}(\cdot \mid c)$.
Varying $c$ therefore varies the member of the family $\mathcal{Q}_{\theta}$ being queried.
But the map $\theta \mapsto p_{\theta}$ is unchanged, since no optimization step or parameter update occurs.
Thus prompting can reveal or organize capabilities already encoded by the pretrained parameters, but it does not itself move the model to a new point in parameter space.
\qed

The proposition is elementary, but it clarifies a major conceptual distinction.
When a prompted model performs a new task, two different explanations are possible.
One is that the pretrained model already contains a sufficiently rich conditional mechanism, and the prompt selects the relevant behavior.
The other is that the model is being genuinely improved by additional parameter adaptation.
Prompting belongs to the first category; instruction tuning belongs to the second.

\subsection*{Instruction tuning and post-pretraining alignment}

Instruction tuning is a parameter update stage applied after broad pretraining.
Suppose one has instruction--response pairs $(u,r)$ drawn from a distribution $P_{\mathrm{inst}}$.
Starting from pretrained parameters $\theta_0$, one optimizes
\[
\mathcal{L}_{\mathrm{inst}}(\theta)
=
\mathbb{E}_{(u,r)\sim P_{\mathrm{inst}}}
\bigl[-\log p_{\theta}(r \mid u)\bigr].
\]
This does not merely select among existing conditionals.
It changes the parameters so that the model becomes easier to steer by instruction-like inputs.

In the current literature, this stage is often grouped under broader labels such as alignment or post-training.
The exact procedures vary: supervised instruction tuning, preference-based optimization, rejection sampling from stronger models, and combinations with retrieval or tool use all appear in practice.
What matters conceptually is the distinction between two questions:
\begin{itemize}
    \item how the broad internal representation was acquired during pretraining,
    \item how the deployed behavior was later shaped toward task following, helpfulness, safety constraints, or conversational format.
\end{itemize}

It is therefore misleading to say that prompting, instruction tuning, and alignment are interchangeable.
Prompting changes the input to a fixed model.
Instruction tuning changes the model itself.
Alignment, in a broad engineering sense, usually refers to post-pretraining procedures intended to shape deployed behavior, of which instruction tuning is one important example.

\subsection*{In-context behavior: what can be stated cleanly}

A striking empirical feature of large pretrained sequence models is that they may behave as if they had adapted to a task merely from examples given in the context.
For instance, one may provide several input--output pairs in a prompt and then ask the model to continue with the output for a new input.
This is commonly called \emph{in-context learning} or \emph{in-context behavior}.

A careful exposition should separate the mathematical fact from the stronger interpretation.
The mathematical fact is simple: because the model computes $p_{\theta}(y\mid c)$ for arbitrary admissible context $c$, one may include demonstrations inside $c$.
The stronger interpretation---that the model is internally implementing a fast learning algorithm over the prompt---is plausible in some cases and supported by several theoretical toy models, but it is not a general theorem about deployed systems.

\begin{definition}[In-context behavior]
Let $\theta$ be fixed.
We say that a model exhibits \emph{in-context behavior} for a task family if its conditional predictor
\[
y \mapsto p_{\theta}(y \mid c)
\]
changes systematically when the context $c$ is augmented by task descriptions or demonstrations, so that performance on the implied task improves without an explicit weight update.
\end{definition}

This definition is intentionally behavioral.
It does not claim that the model literally updates an internal parameter vector in the way an external training algorithm would.
It only claims that the input--output map depends on the context in a task-relevant way.
That is the phenomenon one can observe directly.

\begin{example}[Prompting versus instruction tuning]
Suppose a pretrained language model is asked to classify sentiment.
With no adaptation, one may prompt it by writing:
\begin{quote}
Text: ``The film was unexpectedly moving.''\newline
Label:
\end{quote}
and sampling or scoring candidate outputs such as ``positive'' and ``negative.''
If one instead provides several labeled examples before the query, this becomes few-shot prompting.
By contrast, instruction tuning would optimize the model parameters on many such instruction--response pairs before deployment.
The first is conditional use of a fixed model; the second is parameter adaptation.
\end{example}

\subsection*{Learning at scale: what is algorithmic fact and what is empirical observation}

Foundation models are strongly associated with scale.
To discuss this responsibly, it helps to separate three kinds of statement.

First, there are \emph{architectural and algorithmic facts}.
For a dense transformer with depth $L$, width $d$, and context length $T$, one can analyze parameter count, memory footprint, and per-step computational cost.
At a high level, the parameter count of the feedforward and attention projections scales on the order of
\[
O(Ld^2),
\]
while dense self-attention contributes sequence-length dependence on the order of
\[
O(LT^2d)
\]
in time and roughly $O(T^2)$ per layer in attention-score storage before implementation details are taken into account.
These are not empirical laws but consequences of the architecture.
They explain why larger width, depth, and context require substantial engineering resources.

Second, there are \emph{empirical scaling observations}.
Across substantial but finite training regimes, several studies report that validation loss or perplexity often follows approximate power-law trends as model size, data, and compute increase.
Symbolically, one often writes a relation of the form
\[
\mathcal{L}(N,D,C)
\approx
A N^{-\alpha} + B D^{-\beta} + \cdots,
\]
where $N$ denotes parameter count, $D$ effective training data, and $C$ compute budget.
This notation is useful, but it should not be mistaken for a theorem of learning theory.
The exponents are fitted in restricted regimes, depend on modeling choices, and need not be universal.

Third, there are \emph{design practices and frontier claims}.
Claims about sudden capability transitions, the best post-training recipe, the correct balance between synthetic and human data, or the ideal mixture of retrieval, tools, and prompting are largely empirical and engineering-driven.
Some are well supported by repeated practice.
Very few are settled mathematically.

This separation is essential.
It is true, in a precise computational sense, that larger models usually cost more to train and serve.
It is true, in a repeated empirical sense, that performance often improves predictably over useful ranges as scale increases.
It is not true that current theory provides a complete first-principles account of why broad competence emerges at particular scales or why one post-training pipeline dominates another.

\begin{figure}[h]
\centering
\fbox{%
\begin{minipage}{0.94\textwidth}
\centering
\textbf{Lifecycle of a modern foundation-model pipeline}

\vspace{0.6em}
\begin{tabular}{p{0.82\textwidth}}
\centering Broad data source $P_{\mathrm{pre}}$ \\
\centering $\downarrow$ \\
\centering Pretraining objective $\mathcal{L}_{\mathrm{pre}}$ \\
\centering $\downarrow$ \\
\centering Shared base model $F_{\theta_0}$ \\
\centering $\downarrow$ \\
\centering Post-pretraining shaping: instruction tuning, preference optimization, safety constraints \\
\centering $\downarrow$ \\
\centering Deployment-time interfaces: prompting, retrieval, tools, API calls \\
\centering $\downarrow$ \\
\centering Optional task adaptation: probe, adapter, low-rank update, or full fine-tuning \\
\centering $\downarrow$ \\
\centering Deployed system for downstream use
\end{tabular}
\end{minipage}}
\caption{A print-friendly schematic of the lifecycle.
Pretraining, post-training alignment, prompting, and task adaptation occupy different levels of the pipeline.
The same base model may support several deployment routes, and the mathematical role of each route is different.}
\end{figure}

\subsection*{Why scale matters for representation learning}

Why should large data and large models help at all?
At the chapter level, the most conservative answer is representational.
A broad pretraining distribution exposes the model to many contexts, tasks, styles, and regularities.
A sufficiently expressive architecture can then allocate parameters to internal variables that are not tied to one benchmark label space.
Scale increases the opportunity to learn stable reusable coordinates rather than brittle task-specific shortcuts.

This statement should also be interpreted carefully.
Scale does not guarantee semantic understanding, factual reliability, or robust reasoning in every setting.
It changes the regime in which representation learning takes place.
The resulting models can then often support a larger collection of downstream behaviors through the interfaces described above.
That is already a substantial claim, and it is the one most closely aligned with the rest of this chapter.

\subsection*{What is understood, what is empirical, and what remains design practice}

We can now summarize the epistemic status of the main claims in this area.

\paragraph{Relatively well understood.}
The probabilistic semantics of autoregressive and masked objectives are clear.
The distinction between prompting and parameter adaptation is clear.
The computational scaling of dense architectures in parameters, memory, and sequence length is clear.
The existence of transfer through shared representations is supported both by theory in simplified settings and by overwhelming empirical evidence.

\paragraph{Strongly empirical but stable enough to teach.}
Large-scale pretraining often produces features that transfer broadly.
Instruction tuning often improves task following and interface usability.
 Approximate scaling-law regularities are often visible over substantial ranges.
Prompt design, retrieval augmentation, and post-training methods materially affect deployment behavior.
These statements are not universal theorems, but they are robust enough to form part of modern practice.

\paragraph{Still largely design practice or open theory.}
There is no complete theory of in-context learning in realistic models.
There is no settled first-principles account of why some capabilities appear gradually while others seem to improve abruptly under certain evaluation protocols.
There is no single mathematically justified recipe for post-training alignment.
Questions of robustness, calibration, tool use, long-horizon reasoning, and multimodal grounding remain only partially understood.

That balance matters pedagogically.
A mathematically serious treatment should neither understate what is genuinely established nor overstate what remains heuristic.
Foundation models are important partly because they work, but academic exposition should still distinguish clean formal structure from well-supported empirical regularity and from frontier design judgment.

\subsection*{What to remember}

\begin{itemize}
    \item A foundation model is defined here by its role: broad pretraining, reusable internal structure, and support for many downstream tasks through stable interfaces.
    \item Pretraining objective, adaptation mechanism, and deployment interface are distinct objects and should not be conflated.
    \item Prompting changes context for a fixed model, whereas instruction tuning changes parameters after pretraining.
    \item In-context behavior is an empirical phenomenon of task-sensitive conditioning without explicit weight updates.
    \item Computational scaling of dense architectures is analyzable, but observed scaling laws and post-training recipes are primarily empirical.
\end{itemize}

\subsection*{Common confusion}

\begin{itemize}
    \item \textbf{``Foundation model'' means only ``very large model.''} Size matters in practice, but the more important idea is broad reusable pretraining plus many downstream uses.
    \item \textbf{Prompting and fine-tuning are the same.} They are not: prompting leaves parameters fixed, while fine-tuning updates them.
    \item \textbf{Instruction tuning is identical to pretraining.} It is usually a later stage with different data and different behavioral goals.
    \item \textbf{Scaling laws are theorems.} They are fitted empirical regularities, not universal derivations from first principles.
    \item \textbf{In-context learning is fully understood.} It is observed and partially modeled, but a complete theory for realistic systems is still missing.
\end{itemize}

\subsection*{Exercises}

\begin{enumerate}
    \item Let a pretrained decoder define $p_{\theta}(x_{1:T}) = \prod_{t=1}^{T} p_{\theta}(x_t \mid x_{<t})$. Show directly that a prompt $c$ induces a conditional predictor $q_{\theta,c}(y)=p_{\theta}(y\mid c)$. Why does this formalize prompting without parameter adaptation?
    \item Give one example each of a pretraining objective, an adaptation mechanism, and a deployment interface for a language model. Explain why confusing any two of them leads to conceptual error.
    \item Compare zero-shot prompting, few-shot prompting, instruction tuning, and full fine-tuning. For each, say what changes mathematically: the input context, the parameter vector, or both.
    \item Suppose a dense transformer has depth $L$, width $d$, and context length $T$. Using the chapter's rough scaling discussion, explain qualitatively how training cost changes when one doubles $d$ versus when one doubles $T$.
    \item A student claims that because prompting can produce different behaviors, it should count as learning in exactly the same sense as gradient-based fine-tuning. Give a careful response using the proposition in this section.
    \item Write a short essay comparing two deployment routes for the same base model: prompt-only use versus adapter-based task adaptation. Discuss flexibility, cost, and control.
\end{enumerate}

\subsection*{Notes and references}

For the phrase \emph{foundation model} and a broad systems-level framing, a standard reference is the survey by Bommasani et al.
Canonical large-scale language-model papers include BERT for masked language modeling, GPT-style autoregressive models, T5 for text-to-text transfer, and later large decoder-only systems such as PaLM and related families.
For empirical scaling-law work, the Kaplan et al.~study is the classic early reference, while Hoffmann et al.~(``Chinchilla'') is central for the compute--data tradeoff discussion.
For multimodal large-scale pretraining, CLIP is a standard reference on contrastive vision--language alignment, and Flamingo-style work is representative of multimodal prompting with frozen or partially frozen backbones.
On instruction tuning and post-training alignment, common reference points include the FLAN line of work, InstructGPT, and subsequent preference-optimization methods.
Theoretical work on in-context learning exists in several simplified regimes, including analyses of transformers as implicit learners in stylized settings, but the field remains much less settled than the empirical literature.
